%% ma: L10-B15-0082
%% lop: 10
%% chuong: 6
%% bai: 15
%% dang: 10.15.1
%% loai: DS
%% muc: [NB, TH, TH, VD]
%% dap_an: "SĐĐĐ"
%% nguon: "Ảnh giáo viên gửi 25/11/2026: Bài tập trắc nghiệm – Toán 10 Bài 15 Hàm số (3 phần), Phần II Câu 4"
%% kien_thuc: "Lớp 10 Bài 15 (tập xác định của hàm số); bất phương trình bậc nhất, giao các khoảng (Lớp 10 Bài 2)"
%% trang_thai: chinh-thuc
\begin{ex}
Cho hai hàm số $f(x)=\sqrt{5+x}+\sqrt{5-x}$; $g(x)=\dfrac{3x+4}{\sqrt{x+4}}$. Xét tính đúng sai của các khẳng định sau:
\choiceTF
{Hàm số $g(x)=\dfrac{3x+4}{\sqrt{x+4}}$ xác định khi và chỉ khi $x+4\ge0$.}
{\True Tập xác định của hàm số $f(x)=\sqrt{5+x}+\sqrt{5-x}$ là $D_1=[-5;5]$.}
{\True Tập xác định của hàm số $g(x)=\dfrac{3x+4}{\sqrt{x+4}}$ là $D_2=(-4;+\infty)$.}
{\True Gọi $D_1;D_2$ lần lượt là tập xác định của các hàm số $f(x);g(x)$. Khi đó $(D_1\cap D_2)\setminus(1;6)=(-4;1]$.}
\loigiai{a) Cần $x+4>0$: Sai. b) $-5\le x\le5$: Đúng. c) $x>-4$: Đúng. d) $D_1\cap D_2=(-4;5]$, bỏ $(1;6)$ còn $(-4;1]$: Đúng.}
\end{ex}
